\documentclass[reqno]{amsart} 

\usepackage{amsmath}
\usepackage{amssymb}
\usepackage{hyperref}
\usepackage{amscd}
\usepackage[all]{xy}
\CompileMatrices
\input{xypic}

\theoremstyle{plain}%default
\newtheorem{thm}{Theorem}[section]
\newtheorem{prop}[thm]{Proposition}
\newtheorem{conj}[thm]{Conjecture}
\newtheorem{lemma}[thm]{Lemma}
\newtheorem{cor}[thm]{Corollary}
\newtheorem{cor*}[thm]{Corollary}
\newtheorem*{thm*}{Theorem}
\newtheorem*{prop*}{Proposition}
\theoremstyle{definition}
\newtheorem{defn}[thm]{Definition}
\newtheorem{example}[thm]{Example}
\theoremstyle{remark}
\newtheorem{obs}[thm]{Observation}
\newtheorem{remark}{Remark}

%My general math commands
\newcommand{\Z}{\mathbb Z}
\newcommand{\Q}{\mathbb Q}
\newcommand{\R}{\mathbb R}
\newcommand{\F}{\mathbb F}
\newcommand\vv{\,\vert\,}
\newcommand\sub{\subseteq}
\newcommand\mf[1]{\mathfrak #1}
\newcommand\mb[1]{\mathbb #1}
\newcommand\mc[1]{\mathcal #1}
\newcommand{\CP}{\mathbb{CP}}
\newcommand\LL{\Lambda}
\newcommand\GG{\Gamma}
\newcommand\ot[2]{{#1}\otimes{#2}}
\newcommand\rar[1]{\stackrel{#1}{\longrightarrow}}
\newcommand\dlim{\lim\limits_{\longrightarrow}}
\newcommand\hhom[3]{\operatorname{Hom}_{#1}(#2,#3)}
\newcommand{\rad}{\operatorname{rad}}
\newcommand\dg{\text{deg}\,}
\newcommand\dph{\text{depth}\,}
\newcommand\as{\text{Ass}\,}
\newcommand\rk{\text{rk }}
\newcommand\an{\text{ann}\,}
\newcommand\p[1]{\frak p_{#1}}
\newcommand\fg{finitely generated }
\newcommand\hc[2]{\text{hocolim}_{#1}{#2}}
\newcommand\res[2]{\text{res}^{#1}_{#2}}


\title{Associated primes in group cohomology and a strange result} 
\author{James A. Schafer}

\date{\today}

\address{\begin{flushleft}\quad Department of Mathematics \\
\quad University of Maryland \\
\quad College Park, Maryland 20742
\end{flushleft}}

\email{jas@umd.edu}

\begin{document}
\maketitle

\begin{abstract} In 2010.  Okuyama \cite{O} proved the following.
 
\begin{thm*} Let $G$ be a finite group and $k$ an algebraically closed field of characteristic $p$ dividing the order of $G$. 
\begin{enumerate}\item Let $\frak p$ be an associated prime of $H^*(G,k)$ of $\dim=s$.  Then there exists an elementary abelian $p$-subgroup $E\sub G$ of rank $s$ with $\frak p=\frak p_E$ and depth$\,H^*(C_GE,k)=s$.
\item Let $E$ be an elementary abelian $p$-subgroup of $G$ of rank $s$ and suppose depth$\,H^*(C_GE,k)=s$.  Then $\frak p_E$ is an associated prime of $H^*(G,k)$.
\end{enumerate}\end{thm*}
 
It is shown this result is true for all fields of characteristic $p$. We also show the following result: 
\begin{thm*} If $E$ is an elementary abelian $p$-subgroup of $G$ and $\dph H^*(G,k)=\rk E$ then up to conjugation $E$ contains the maximal elementary abelian subgroup contained in the center of a Sylow $p$-subgroup of $G$. \end{thm*}

\end{abstract}

\section{Okuyama's result}
\begin{thm*} Let $G$ be a finite group and $k$ a field of characteristic $p$ dividing the order of $G$. 
\begin{enumerate}\item Let $\frak p$ be an associated prime of $H^*(G,k)$ of $\dim=s$.  Then there exists an elementary abelian $p$-subgroup $E\sub G$ of rank $s$ with $\frak p=\frak p_E$ and depth$\,H^*(C_GE,k)=s$.
\item Let $E$ be an elementary abelian $p$-subgroup of $G$ of rank $s$ and suppose depth$\,H^*(C_GE,k)=s$.  Then $\frak p_E$ is an associated prime of $H^*(G,k)$.
\end{enumerate}\end{thm*}
  
If one wishes, by using results of Serre and Wilkerson \cite{W} the first part of the theorem in the general case is easy to show.  Their results show that if $\frak p$ is an associated prime of $H^*(G,k)$ then $\frak p=\frak p_E$ for some elementary abelian $p$-subgroup of $G$ with $\rk E=\dim \frak p$.  By Quillen \cite{Q} there exists $\tau_E\in \frak p_E$ and $\tau_E\notin \frak p_{E'}$ for all $E'$ with $\rk E'=\rk E$. If $\frak p_E=\an\,y$ then $\tau_E\,y=0$ and therefore $\res G{E'}\tau_E\,y=0$.
Since $\res G{E'}\tau_E$ is not in the radical of $H^*(E')$ it is a non-zero divisor and therefore a non-zero divisor in $H^*(C_GE')$.  Hence $\res G{C_GE'}\,y=0$. Now a result of Carlson \cite{C} says $\bar y=\res G{C_GF}y\neq 0$ for some elementary abelian $F$ with $\rk F=\rk E$. Hence $\bar y=\res G{C_GE}\,y\neq 0$. Let $J=\an \bar y$. $\res G{C_GE}(\frak p_E)\sub J$ 
$$\rk E=\dim\frak p_E\geq\dim\,(\res G{C_GE})^{-1}J=\dim J \geq\dim Q\geq\dph H^*(C_GE)\geq \rk E$$
where $Q$ is an associated prime of $H^*(C_GE)$. 

Okuyama's proof in both directions depends on results of Green \cite{G} which are only valid for $p$-groups and a result of Benson \cite{B1} whose proof is questionable.  I asked AI [Claude Opus5.5 (Antropic 2026)] to give a proof of Green's results valid for arbitrary groups and find a proof of Benson's result and then write a paper doing so. The resulting paper is available at http://www.terpconnect.umd.edu/$\sim$jas.  It shows that Okuyama's theorem only needs $k$ to be infinite and not algebraically closed.  This is because of Lemma 6.3 in the paper, which definitely uses $k$ infinite in the proof. Replacing that proof with the proof of the lemma below, valid for any field, allows one to conclude Okuyama's theorem is valid for all fields.

\begin{lemma} Let $Q$ be a finite group acting effectively and linearly on a finite dimensional $k$-vector space $U$ and hence on the symmetric algebra $S(U)$. Then there is $f\in S(U)$ such that $\{{}^qf\,\vert\, q\in Q\}$ is a linearly independent set.\end{lemma}
\begin{proof} Identify $S(U)$ with $k[U^*]$.  The effective action of $Q$ induces an effective action of $Q$ on $K=k(U^*)$. Hence $K$ is a Galois extension of $K^Q$ with Galois group $Q$.  By the Normal Basis theorem there exists $\alpha\in K$ with $\{{}^q \alpha\}$ a $K^Q$-basis of $K$ and therefore linearly independent over $k$.  If $\alpha=\beta/\gamma$ with $\beta,\,\gamma\neq 0$ in $S[U]$, let $\rho=\prod{}^q\gamma\in K^Q$.  Then if  
$\bar\alpha=\alpha\,\rho\in S[U]$, $\{{}^q\bar\alpha\}=\{\rho\,{}^q\alpha\}$ is a linearly independent set in $S(U)$.\end{proof}

\section{A strange result}

Let $G$ be a finite group and $k$ a field of characteristic $p$ dividing the order of $G$.  Let $E_0$ be the unique (up to conjugation) maximal elementary abelian $p$-subgroup of $G$ contained in the center of a Sylow $p$-subgroup of $G$.  Let $C_GE$ denote the centralizer of $E$ in $G$.

\begin{thm*} If $E$ is elementary abelian $p$-subgroup of $G$ and $\dph H^*(C_GE)=\rk E$, then $E_0\sub E$ up to conjugation.\end{thm*}
\begin{proof} By \cite{O}, $\frak p_E$ is an associated prime of $H^*(G)$ and by \ref{fine}, $E$ must contain $E_0$.
\end{proof}

Why is this result strange?  What has the depth of $H^*(C_GE,k)$ have to do with the placement of $E$ in the lattice of elementary subgroups of $G$?

\begin{prop} Let $f:A\to B$ be a ring homomorphism. If $\mf q\sub B$ is a primary ideal with $\mf p=\rad\,\mf q$, then $f^{-1}\mf q\sub A$ is a primary ideal of with $\rad\,f^{-1}\mf q=f^{-1}\mf p$.
\end{prop} 
\begin{proof} If $f(uv)\in\mf q$ with $f(v)\notin\mf q$, then $f(u)^s=f(u^s)\in\mf q$. 
$$ y^s\in f^{-1}\mf q\iff f(y)^s=f(y^s)\in\mf q\iff y\in f^{-1}\rad\mf q.$$\end{proof}

\begin{prop} Let $A\sub B$ be Noetherian.  Then $\as A\sub\{\mf p\bigcap A\,\vert\, \mf p\in\as B\}$.
\end{prop}
\begin{proof}
Since $A$ and $B$ are Noetherian, if $(0)=\bigcap\mf q$ is a minimal primary decomposition of $(0)\sub A$ then $\as A=\{\mf p=\rad q\}$ and similarly for $B$.  Therefore if $(0)=\bigcap\mf q$ is a minimal primary decomposition of $(0)\sub B$ with primes $\{\mf p\}$ then $(0)= \bigcap \mf q\cap A$ is a primary decomposition of $(0)\sub A$ with primes $\{\mf p\cap A\}$ by the above proposition but not necessarily minimal. However it can always be made minimal by a) amalgamating primary ideals with then same radical and b) eleminating a primary ideal $\frak q$ with $\bigcap_{\frak q'\neq\frak q}\sub\frak q$ and therefore $\as A\sub\{\mf p\bigcap A \,\vert\, \mf p\in\as B\}$.

\end{proof}

\begin{thm}\cite{Q} $\frak p_E\sub \frak p_{E'}$ iff $E'$ is conjugate to a subgroup of $E$. \end{thm}

\begin{cor}\ Let \,$E$ be elementary $p$-abelian.  Then there exists $\tau_E\in\frak p_E$ such that $\tau_{E}\notin\frak p_{E'}$ for any elementary abelian $E'$ not conjugate to a subgroup of $E$.\end{cor}
\begin{proof} Let $W=\{E'\in\mathcal A(G)\vv E'\,\text{not conjugate to a subgroup of}\,E\}$. Then $\frak p_E\not\subseteq \frak p_{E'}$ for any $E'\in W$.  By Prime Avoidence $\frak p_E\not\subseteq\bigcup_W\frak p_{E'}$.  Therefore there exists $\tau_E\in\frak p_E$ such that $\tau_{E}\notin\frak p_{E'}$ for any $E'\in W$. \end{proof} 

\begin{remark} Since if $g \in G$ then $\frak p^G_E=\frak p^G_{{}^gE}$ we may always assume the elementary abelian $p$-subgroup in $\frak p^G_E$ has $E\sub P$.
\end{remark}

\begin{cor}\label{fine} Let $E_0$ be the maximal $p$-elementary abelian subgroup of the center of $P$ where $P$ is a Sylow $p$-subgroup of $G$.  Then if $\frak p_E\in\as H^*(G,k)$, $E_0\sub E$.\end{cor}
\begin{proof} Since $H^*(G,k)\sub H^*(P,k)$ the above proposition gives $\as(H^*(G,k))\sub\{\frak p^P_F\cap H^*(G,k)\vv \frak p^P_F\in\as H^*(P,k)\}$.  Hence 
$$\frak p^G_E=\frak p^P_F\cap H^*(G)=\frak p^G_F, \text{ with } \frak p^P_F\sub \as\,H^*(P,k).$$
That is we may always assume that up to conjugation $\frak p^P_E\in\as H^*(P,k)$.  

Let $x\in Z(P)$ be an central element of order $p$ of $P$ and let 
$$\frak p^P_E=\ker\{H^*(P,k)\to H^*(E,k)/\text{Rad}\}\in\as H^*(P,k).$$  Suppose $x\notin E$. Then $E'=(x)$ is not conjugate in $P$ to a subgroup of $E$ and so there exists $\tau\in\frak p_E\sub H^*(P,k)$ with $\tau\notin\frak p^P_{E'}$.  Therefore 
$$0\neq\pi\circ\text{Res}^P_{E'}(\tau)\in H^*(E',k)/\text{Rad}. $$  That is $\text{Res}^P_{E'}(\tau)\in H^*(E',k)$ is not nilpotent.  Since $H^*(E',k)=P[\beta\,\bold u]\otimes \Lambda[\bold u]$ this implies $\text{Res}^P_{E'}(\tau)$ is a non-zero-divisor. Since $E'$ is central in $P$, a result of Duflot, \cite{D}, says $\tau\in H^*(P,k)$ is a non zero-divisor and since $\frak p^P_E=\text{ann}_{H^*(P,k)}(y)$ for some $y\in H^*(P,k)$, $\tau\notin\frak p^P_E$. But $\tau\in \frak p_E$ and therefore every central element of $P$ is contained in $E$ and $E_0\sub E$.
\end{proof}

\begin{thebibliography}{B1}
\bibitem{B} Benson, D.J., \emph{ Dickson invariants, regularity and computation in group cohomology}, Illinois J. Math., \textbf{48}, (2004), 171-197.
\bibitem{B1} Benson, D.J., \emph{The image of the transfer map} Arch. Math. (Basel) \textbf{61}, (1993), 7-11.
\bibitem{C} Carlson, J.F., \emph{Depth and transfer in the cohomology of groups}, Math. Zeit., \textbf{218}, (1995), 461-468.
\bibitem{D} Duflot, J., \emph{Depth and equivariant cohomology}, Comment. Mah. Helv., \textbf{56},(1981),671-637.
\bibitem{G} Green, D.J., \emph{On Carlson's depth conjecture in group cohomology}, Math.\ Z.\ \textbf{244} (2003), 711--723; arXiv:math/0206127.
%\bibitem{OS} Okuyama, T. and Sasaki, H., \emph{Homogeneous systems of parameters in cohomolgy algebras of finite groups}, Archiv der Math. \textbf{82} (2004), 110-121.
\bibitem{O} Okuyama, T. \emph{A remark on associated primes in the cohomology of a finite group}, RIMS K\^oky\^uroku, \textbf {1679}, (2010), 113-117.
\bibitem{Q} Quillen, D. \emph{The spectrum of an equivariant cohomology ring I, II}, Ann.\ of Math.\ \textbf{94} (1971), 549--572 and 573--602.
\bibitem{W} Wilkerson, C.W., \emph{A primer on the Dickson invariants}, Proc. of the Northwestern Homotopy Theory Conf., Contemp. Math., \textbf{19}, AMS, (1983), 421-434. 
\end{thebibliography}


\end{document}
\section*{Associated Primes in \fg graded $k$-algebras}

\begin{lemma}\label{two} Let $B\sub A$ be Noetherian rings with $A$ \fg as an $B$-module. Let $M$ be a \fg $A$-module and $Q$ a primary $A$-submodule of $M$.  Then $Q$ is an $B$-primary submodule of $M$ and if $P=r(\text{ann}_A(M/Q)$ is the associated prime then the associated prime for $Q$ as a primary $B$-module is $P\bigcap B=r(\text{ann}_B(M/Q)$.\end{lemma}
\begin{proof} The first part is obvious and the second part is immediate since $\text{ann}_B(L)=B\bigcap\text{ann}_S(L)$ for any $A$-module $L$.\end{proof}

\begin{prop}\label{three} If $B\sub A$ are Noetherian rings with $A$ \fg as an $B$-module. Let $M$ be a \fg $A$-module. If $\as_AM=\{P_1,\dots,P_s\}$ then $\as_BM=\{P_1\cap B,\dots,P_s\cap B\}$. In particular $\omega_BM=\omega_AM$.\end{prop}
\begin{proof} If $(0)=\bigcap Q_i$ is a minimal primary decomposition of $(0)$ by primary $A$-submodules of $M$ with corresponding primes $\{P_i\}$ then it is also a $B$-primary decomposition with primes $\{P_i\bigcap B\}$ by the lemma.  This decompostition may not be minimal. 	If $\bar Q_{r_j}=\bigcap\{Q_i\vv P_i\bigcap B=P_{r_j}\bigcap B\}$ then $(0)=\bigcap\bar Q_{r_j}$ gives an $B$-primary decomposition with distinct primes. It must be minimal for if $\bigcap_{j'\ne j}\bar Q_{r_{j'}}\sub \bar Q_{r_j}$ then $\bigcap_{s\ne r_j}Q_s\sub Q_{r_j}$ which is impossible since $(0)=\bigcap Q_i$ is a minimal primary decomposition.\end{proof}

We now investigate how minimal primary decompositons act under submodules.  We assume $A$ is a commutative (or skew-commutative), Noetherian ring, $M$ is a \fg module and $N\sub M$ is a submodule.

\begin{prop}\label{four} Let $M$ be a \fg module and $N\sub M$. Let $Q\sub M$ be a primary submodule of $M$ with prime $P=r(\text{ann}_AM/Q)$. If $N\cap Q\neq N$ then $N\cap Q$ is a primary submodule of $N$ with prime $r(\text{ann}_AN/(N\bigcap Q))$ equal to $P$. \end{prop}
\begin{proof} It is sufficient to show $\as(N/(N\cap Q)=\{P\}$.  But 
$$N/(N\cap Q)\simeq (N+Q)/Q\sub M/Q.  $$ Hence $\emptyset\neq\as(N/(N\cap Q))\sub\as(M/Q)=\{P\}$ and we are done.\end{proof}

\begin{thm}\label{non} Let $B$ be a subring of $A$ and $M$ a \fg $A$-module and $N\sub M$. Suppose $\as_AM=\{P_1,\dots,P_s\}$.  Then
\begin{enumerate}\item There exists a non-necessarily minimal $B$-primary decomposition $(0)=\bigcap\,Q\cap N\sub N$, where $Q$ are $A$-submodules of $M$.
\item If $\bar P=\text{ann}_B(N/(Q\cap N)$ is the prime associated to $Q\cap N$, then $\bar P\in\{P\cap B\vv P\in\as A\}$.
\end{enumerate}\end{thm}
\begin{proof} Let $(0)=\bigcap Q$ be a minimal $A$-primary decomposition of $(0)\sub M$ and $N\sub M$ a $B$-submodule.  If $P_Q$ is the associated prime of $Q$ and $Q_P=\bigcap\{Q'\vv P_{Q'}\cap B=P_Q\cap B\}$ then $(0)=\bigcap Q_P$ is a minimal $B$-primary decomposition of $M$ by $A$-submodules with primes $\{P\cap B\vv P\in\as_AM\}$ by \ref{three}.

By \ref{four} $(0)=\bigcap\{Q_P\cap N\vv Q_P\cap N\neq N\}$ is a not necessarily minimal $B$-primary decomposition of $(0)$ with primes contained in the set $\{P\cap B\}$.  %
%If $B\sub Q_P$ then $1_B=1_A\in Q_P$ which implies $A=Q_P$ since $Q_P$ is an ideal in $A$.  Therefore by \ref{four} $(0)=\cap_P\,(Q_P\cap B)$ is a not necessarily minimal $B$-primary decompostion of $(0)\sub B$. Moreover since $r(Q\cap B)=r(Q)\cap B$, $Q\cap B$ is $P\cap B$-primary.
\end{proof}
\begin{remark} If $M=A$ and $N=B$ then $Q\cap B=B$ is imposible since $1_B=1_A\notin Q$ and $Q$ is an $A$-module. Hence in the non-minimal primary decomposition of $(0)\sub B$ the set of primes associated to $\{Q\cap B\}$ equals $\{P\cap B\vv P\in\as_A(A)\}$. \end{remark}

\begin{thm}\cite{Q} $\frak p_E\sub \frak p_{E'}$ iff $E'$ is conjugate to a subgroup of $E$. \end{thm}

\begin{cor}\ Let \,$E$ be elementary $p$-abelian.  Then there exists $\tau_E\in\frak p_E$ such that $\tau_{E}\notin\frak p_{E'}$ for any elementary abelian $E'$ not conjugate to a subgroup of $E$.\end{cor}
\begin{proof} Let $W=\{E'\in\mathcal A(G)\vv E'\,\text{not conjugate to a subgroup of}\,E\}$. Then $\frak p_E\not\subseteq \frak p_{E'}$ for any $E'\in W$.  By Prime Avoidence $\frak p_E\not\subseteq\bigcup_W\frak p_{E'}$.  Therefore there exists $\tau_E\in\frak p_E$ such that $\tau_{E}\notin\frak p_{E'}$ for any $E'\in W$. \end{proof} 

\begin{remark} Since if $g \in G$ then $\frak p^G_E=\frak p^G_{{}^gE}$ we may always assume the elementary abelian $p$-subgroup in $\frak p^G_E$ has $E\sub P$.
\end{remark}

\begin{cor}\label{fine} Let $E_0$ be the maximal $p$-elementary abelian subgroup of the center of $P$ where $P$ is a Sylow $p$-subgroup of $G$.  Then if $\frak p_E\in\as H^*(G,k)$, $E_0\sub E$.\end{cor}
\begin{proof} By the remark after Theorem 0.6 the associated primes of $H^*(G,k)\sub H^*(P,k)$ are a subset of $\{\frak p^P_E\cap H^*(G,k)\vv \frak p^P_E\in\as H^*(P,k)\}$.  Hence 
$$\frak p^G_E=\frak p^P_{\tilde{E}}\cap H^*(G)=\frak p^G_{\tilde{E}}, \text{ with } \frak p^P_{\tilde{E}}\sub \as\,H^*(P,k).$$
That is we may always assume that up to conjugation $\frak p^P_E\in\as H^*(P,k)$.  

Let $x\in Z(P)$ be an central element of order $p$ of $P$ and let 
$$\frak p^P_E=\ker\{H^*(P,k)\to H^*(E,k)/\text{Rad}\}\in\as H^*(P,k).$$  Suppose $x\notin E$. Then $E'=(x)$ is not conjugate in $P$ to a subgroup of $E$ and so there exists $\tau\in\frak p_E\sub H^*(P,k)$ with $\tau\notin\frak p^P_{E'}$.  Therefore 
$$0\neq\pi\circ\text{Res}^P_{E'}(\tau)\in H^*(E',k)/\text{Rad}. $$  That is $\text{Res}^P_{E'}(\tau)\in H^*(E',k)$ is not nilpotent.  Since $H^*(E',k)=P[\beta\,\bold u]\otimes \Lambda[\bold u]$ this implies $\text{Res}^P_{E'}(\tau)$ is a non-zero-divisor. Since $E'$ is central in $P$, a result of Duflot, \cite{D}, says $\tau\in H^*(P,k)$ is a non zero-divisor and since $\frak p^P_E=\text{ann}_{H^*(P,k)}(y)$ for some $y\in H^*(P,k)$, $\tau\notin\frak p^P_E$. But $\tau\in \frak p_E$ and therefore every central element of $P$ is contained in $E$ and $E_0\sub E$.

\end{proof}

\begin{thebibliography}{B1}
\bibitem{B} Benson, D.J., \emph{ Dickson invariants, regularity and computation in group cohomology}, Illinois J. Math., \textbf{48}, (2004), 171-197.
\bibitem{B1} Benson, D.J., \emph{The image of the transfer map} Arch. Math. (Basel) \textbf{61}, (1993), 7-11.
\bibitem{C} Carlson, J.F., \emph{Depth and transfer in the cohomology of groups}, Math. Zeit., \textbf{218}, (1995), 461-468.
\bibitem{D} Duflot, J., \emph{Depth and equivariant cohomology}, Comment. Mah. Helv., \textbf{56},(1981),671-637.
\bibitem{G} Green, D.J., \emph{On Carlson's depth conjecture in group cohomology}, Math.\ Z.\ \textbf{244} (2003), 711--723; arXiv:math/0206127.
%\bibitem{OS} Okuyama, T. and Sasaki, H., \emph{Homogeneous systems of parameters in cohomolgy algebras of finite groups}, Archiv der Math. \textbf{82} (2004), 110-121.
\bibitem{O} Okuyama, T. \emph{A remark on associated primes in the cohomology of a finite group}, RIMS K\^oky\^uroku, \textbf {1679}, (2010), 113-117.
\bibitem{Q} Quillen, D. \emph{The spectrum of an equivariant cohomology ring I, II}, Ann.\ of Math.\ \textbf{94} (1971), 549--572 and 573--602.
\bibitem{W} Wilkerson, C.W., \emph{A primer on the Dickson invariants}, Proc. of the Northwestern Homotopy Theory Conf., Contemp. Math., \textbf{19}, AMS, (1983), 421-434. 
\end{thebibliography}

\end{document}
